\chapter{Ecoulements incompressibles et irrotationnels non visqueux}

Si la masse volumique $\rho$ est une constante du mouvement, alors par l'équation de continuité \eqref{eq:cons_continuite_lineique}, le champ de vitesse est non-divergent

\begin{equation}
  \Div \vec{u} = 0\text{ pour un écoulement incompressible}
\end{equation}


\section{Théorème de Bernouilli}

Plaçons nous dans les conditions du chapitre \ref{ch:euler}: écoulement laminaire, non-visqueux, dans lequel la seule force interne est la force de pression telle que nous l'avons définie pour établir l'équation d'Euler. 
\begin{equation}
  \rho \, \Dt{\vec{u}} = - \rho \, \grad \gpot - \grad p.
\end{equation}

Nous utilisons deux hypothèses supplémentaires: l'écoulement est stationnaire ($\pdt{}=0$) et incompressible. On utilise par ailleurs la forme alternative du tenseur d'advection (eq: \ref{eq:euler_conserv}):

\begin{equation}
  \rho (\curl \u)\cross u + \grad(\frac{1}{2}\u^2) = - \rho \, \grad \gpot - \grad p.
\label{eq:Bernouilli0}
\end{equation}

Comme $\rho$ est constant et non-nul, nous pouvons nous autoriser à diviser tous les termes par $\rho$ et l'introduire sous la dérivée spatiale: 
\begin{equation}
  (\curl \u)\cross u = - \grad (\gpot + \grad \frac{p}{\rho} - \frac{1}{2}\u^2) \eqdef - \grad H
\label{eq:Bernouilli}
\end{equation}

La quantité $H$ a les dimensions d'une énergie massique, et on y reconnaît les termes évoquant une énergie cinétique par unité de masse, le géopotentiel, et la pression. L'ensemble forme une quantité conservée le long du mouvement. En effet: 

\begin{equation}
  (\u \cdot \grad) H = \u \cdot  (\curl \u)\cross u  = 0
\end{equation}

\begin{key}{Théorème de Bernouilli}
  Si un fluide idéal est en écoulement stationnaire, alors $H$ est constant le long d'une ligne de courant (\cite{acheson01aa}, p. 10).
\end{key}

\section{Équation de la vorticité}

Reprenons l'équation de Bernouilli dans sa version non-stationnaire: 

\begin{equation}
 \pdt u +  (\curl \u)\cross u = - \grad H
\end{equation}

Nous définisson la \cindex{vorticité} comme le rotationnel de la vitesse (cf. ch. \ref{ch:euler}, section \ref{para:Vortex}): $\omega \eqdef \curl \u$. En prenant le rotationnel de part et d'autre de l'équation de Bernouilli  on obtient: 

\begin{equation}
 \pdt \vort +  \curl (\vort \cross u) = 0, 
\end{equation}
%
car le rotationnel d'un gradient est toujours nul. 

On utilise les identités connues en calculus pour développer le rotationnel du produit vectoriel: 

\begin{equation}
 \pdt \vort (\u\cdot \grad)\vort - (\vort\cdot\grad)\u + \vort\Div \u - \u \Div \vort
\end{equation}
%
Si l'écoulement est non-divergent, $\Div u=0$ et de plus $\Div \vort$ est toujours nul par construction de la vorticité. On trouve: 

\begin{equation*}
 \pdt \vort (\u\cdot \grad)\vort = \vort \cdot \grad \u
\end{equation*}
%
ou encore
\begin{equation}
 \Dt\vort = \vort \cdot \grad \u
\label{eq:vort}
\end{equation}
%
C'est l'équation de la vorticité écrite ici pour un écoulement incompressible. Il est remarquable que la vorticité évolue, le long de l'écoulement, indépendamment de la pression et du géopotentiel ! 

\todo[inline]{ dire qu'elle exprime la conservation du moment angulaire. Sans moment de force externe, pas de génération de vorticité nette à l'intérieur du fluide. }

Plus fort encore, si l'écoulement est bi-dimensionnel (par exemple, dans le plan $x,y$, alons $\vort = \omega \ez$ et

\begin{equation}
 \Dt\omega = 0. 
\label{eq:vort2d}
\end{equation}


Cette équation justifie l'existence de particules  de vorticité. Les 'petits' vortex sont advectés avec le flux moyen, comme s'il s'agissait de particules dans un champ 2D. 

Si, d'autre part, l'écoulement est tri-dimensionnel, laminaire est stable, avec une vorticité nulle en condition initiale $\vort=0$, aucune vorticité ne sera créée. Reprenons alors l'équation \eqref{eq:Bernouilli0} qui nous a servi à établir le théorème de Bernouilli: 

\begin{displaymath}
  \rho (\curl \u)\cross u + \grad(\frac{1}{2}\u^2) = - \rho \, \grad \gpot - \grad p.
\label{eq:Bernouilli0}
\end{displaymath}

Comme $\vort = \curl \u =0$, on a 

\begin{displaymath}
  \grad(\frac{1}{2}\u^2) = - \grad \gpot - \grad {p}{\rho}.
\end{displaymath}

Ceci nous donne un "Bernouilli plus": si le flot est non-rotationnel, alors la quantité $H$ est constante, non-seulement le long d'une ligne de courant, mais absolument partout dans le fluide. 

\begin{key}{Bernouilli PLUS}
  Dans un fluide parfait, irrotationnel, la quantité $\gpot + \frac{p}{\rho} + \frac{u^2}{2}$ est constante
\end{key}



\begin{application}{Écoulement libre non-turbulent}
Soit une cuve à l'air libre et un orifice au bas de la cuve, également à pression atmosphérique : la vitesse du fluide en sortie est alors directement déterminée par la différence d'énergie potentielle : $\Delta u/2 = \Delta \Phi$. 

\centering
    \captionsetup{type=figure}
    \input{Pdftex/02_ecoulement_libre.pdf_tex}
    % Use \captionof to generate the figure number and caption
    \raisebox{1em}{
    \parbox{5cm}{\caption{Schéma de la cuve et de l'écoulement libre non-turbulent à travers l'orifice.} \label{fig:ecoulement-libre-application}}}
\end{application}

On évoque souvent l'effet Magnus ou encore la portance d'une aile d'avion comme application du théorème de Bernouilli. 

Le raisonnement est qu'à géopotentiel égal, le long d'une ligne de flux, la pression diminue quand la vitesse augmente. Si la vitesse \emph{au dessus} de l'aile d'avion est plus grande qu'en dessous, alors la différence de pression associée au théorème de Bernouilli génère une portance. 

L'effet Magnus porte sur un objet en rotation (par exemple un ballon). Ici encore, si l'on peut prouver que la vitesse du fluide diffère de part et d'autre du ballon, on explique une force résultante. 

Cependant, il n'est pas trivial de prouver cette différence de vitesse. Nous allons esquisser le raisonnement formel et ensuit réexaminer le cas du ballon (d'abord) et de l'aile d'avion ensuite. 

\section{Équation de Laplace complexe et conditions de Cauchy-Rieman}

\subsection{Écoulement potentiel complexe}

Plaçons-nous dans le cas particulier d'un écoulement bi-dimensionel.  Nous avons vu dans la section \ref{sect:courant} que si la vitesse dérive d'une fonction de courant, $\uh = \curl \sfn \ez$, alors l'équation de continuité est automatiquement satisfaite. 

Si, par ailleurs, le flux $\uh$ dérive d'un potentiel scalaire $\phi$ définit tel que: 

\begin{displaymath}
  u = \pd{\vpot}{x}, \quad \text{et}\quad u = \pd{\vpot}{y}, 
\end{displaymath}

il est automatiquement irrotationnel. Ainsi, si un écoulement dérive à la fois d'une fonction de courant et d'un potentiel scalaire, il est à la fois irrotationel \emph{et} incompressible. 

Ainsi, si on applique la condition de non-divergence à l'expression dérivée de la fonction de courant, on a $\lapl\sfn=0$, et la condition de champ irrotationel à l'expréssion dérivée du potentiel scalaire: $\lapl\vpot=0$. Il s'agit de deux équations de Laplace, qui peuvent aussi s'exprimer en coordonnées polaires:

\begin{equation}
\lapl \sfn = \frac{1}{r}\pd{}{r}\left(r\pd{\sfn}{r}\right) + \frac{1}{r^2}\pd[2]{\sfn}{\theta} = 0. 
  \label{eq:laplace_cp}
\end{equation}

Pour qu'un écoulement puisse dériver à la fois d'un potentiel scalaire et d'une fonction de courant, il faut
que les vitesses associées aux deux définitions soient identiques, soit: 

\begin{equation}
  \pd{\vpot}{x} = \pd{\sfn}{y}\quad\text{et}
  \pd{\vpot}{y} = -\pd{\sfn}{x}. 
  \label{eq:cauchy-riemann}
\end{equation}

Ces équations sont l'expression des \cindex{conditions de Cauchy-Riemann}, qui garantissent l'existence d'une dérivée spatiale à la quantité complexe $\vcomp =\vpot + i \sfn$ (cf. e.g. \cite[sect. 6.2, p. 399]{arfken01aa}). 

Ainsi, si l'on exprime la position avec une variable complexe $z = u + i\,v$, une fonction est qualifiée de \cindex{fonction analytique} lorsque sa dérivée spatiale ne dépend pas de l'orientation de l'incrément.  Concrètement: 

\begin{displaymath}
  \begin{split}
    \od{\vcomp}{z} &= \lim_{\delta x \rightarrow 0}\frac{\delta \vcomp}{\delta x} = \pd{\vcomp}{x} = \pd{\vpot + i\,\sfn}{x} \\
                   &= \lim_{\delta y \rightarrow 0}\frac{\delta \vcomp}{\delta y} = \pd{\vcomp}{y} = \pd{\vpot + i\,\sfn}{y}, 
  \end{split}
\end{displaymath}

On vérifie que les deux expressions sont équivalentes et satisfont: 

\begin{displaymath}
    \od{\vcomp}{z} = u - iv
\end{displaymath}

Nous aboutissons donc à un résultat remarquable et fondamental
\begin{key}{Condition suffisante pour un écoulement potentiel}
toute fonction  qui soit au moins deux fois dérivable dans le plan complexe, satisfait les équations de Laplace dans le plan complexe. 
\end{key}

Ce résultat nous permet de construire quelques écoulements types qui vont nous servir pour construire des écoulements satisfaisant à des conditions frontières et de symmétrie spécifiques. 


\subsection{Écoulements complexes potentiels type}

\paragraph{Écoulement constant}

L'écoulement associé à $\vcomp = U\,z$ a pour vitesse $U\ex$. 


\paragraph{Source}

On génère une source ponctuelle d'amplitude $m$ avec le champ potentiel

\begin{displaymath}
  U = \frac{m}{2\pi}\ln z. 
\end{displaymath}

Le logarithme complex $z$ peut s'estimer en écrivant $z$ sous la forme d'un phaseur complexe $z=re^{i\theta}$. On voit alors que 

\begin{displaymath}
  \ln z = \ln (re^{i(\theta + 2\pi k)}) = \underbrace{\ln r}_{\vpot} + i\underbrace{\theta + 2\pi k}_{\sfn}. 
\end{displaymath}

Attention cependant: l'angle n'est défini qu'à un terme $2\pi k$ près, et c'est crucial. 

On vérifie qu'il s'agit d'un champ divergent, symmétrique autour de l'origine, telque pour tout contour autour de l'origine, le flux net s'échappant du contour est $m$. L'écoulement est donc non-divergent et irrotationnel en tout point, sauf à l'origine où une singularité permet la création d'un fluide. 

\begin{figure}[h!]
    \centering
    
    % --- Figure 1 : Source ---
    \begin{minipage}[t]{0.49\textwidth}
        \centering
        \input{Xp/potential_flow_source.eepic}
        \caption{Potentiel d'écoulement généré par une unique source (débit massique positif). Les lignes de courant (streamlines) sont indiquées par des traits pleins et représentent la fonction de courant $\sfn$ (Streamfunction). Les lignes de potentiel de vitesse (équipotentielles) sont tracées en rouge tiretés et représentent le potentiel de vitesse $\vpot$.}
        \label{fig:potential-source}
    \end{minipage}%
    \hfill%
    %
    % --- Figure 2 : Circulation ---
    \begin{minipage}[t]{0.49\textwidth}
        \centering
        \input{Xp/potential_flow_circulation.eepic}
        \caption{Potentiel d'écoulement généré par une unique circulation (vortex ponctuel). Les lignes de courant sont indiquées par des traits pleins ($\sfn$). Les lignes de potentiel de vitesse (équipotentielles) sont tracées en rouge tireté ($\vpot$) et sont ici explicitement étiquetées pour identification, mais exprimées en degrés. .}
        \label{fig:potential-circulation}
    \end{minipage}
\end{figure}


Sur la figure \ref{fig:potential-source}, on repère les lignes de courant en train plain. On se rappelle que le débit volumique est donné par la différence entre deux valeurs de la fonction de courant. C'est ici que le terme $2\pi k$  est important, et crée une ambiguité sur la valeur de la fonction de courant. $k$ doit être choisi de façon à ne pas créer d'anomalie lorsqu'on passe de 359 à 360 degrés. Cela explicite le fait que la fonciton $\ln z$ n'est pas définie de façon unique sur tout le domaine. Ce qui est défini est sa différence entre deux points du domaine (c'est un cas de \cindex{spirale de Riemann}). 

\paragraph{Vortex ponctuel}

Nous avons déjà rencontré le vortex ponctuel, de circulation $\Gamma$, dans la section \label{para:vortex_ponctuel}. Il répond au potentiel complexe 

\begin{displaymath}
  U = \frac{i\Gamma}{2\pi}\ln z. 
\end{displaymath}

Il génère un diagramme similaire à celui de la source ponctuelle (Figure \ref{fig:potential-circulation}, mais les rôles de la fonction de courant et du potentiel de vitesse sont échangés, avec la même attention à apporter à l'ambiguité de la valeur, cette fois, du potentiel de vitesse. 


%
% \begin{displaymath}
% U_\text{dipole} =  \lim_{\varepsilon \rightarrow 0) 
%   \frac{m}{2\pi}\left( \ln(z*\varepsilon) - \ln(z-\varepsilon) \right)
% \end{displaymath}
%
%
% \paragraph{Circulation}
%

Nous allons voir comment utiliser la puissance des potentiels complexes pour trouver des solutions d'écoulements dans des cas particuliers. 

\subsection{Écoulement autour d'un cylindre (sans circulation)}

Considérons un écoulement horizontal laminaire, de vitesse $U$. Nous savons qu'il répond au potentiel complex $U\,z$. 

Nous exigeons maintenant qu'il contourne un cylindre de rayon $R$ placé à l'origine. Le potentiel doit donc être adapté en appliquant un perturbation dont la forme doit répondre à des exigences de symmetries:

\begin{itemize}
  \item elle implique la seule dimension $R$
  \item symmétrique par rapport à l'origine
  \item s'annule loin du centre
\end{itemize}

La perturbation doit par ailleurs être telle que le potentiel complexe s'annule (ou soit constant) à la frontière du cylindre. Cela va générer par construction un ligne de courant tout au long du cylindre, qui interdit tout flux à travers la paroi du cylindre. 

La fonction $\vcomp(z) = U ( z + \frac{R^2}{z})$ respecte toutes ces conditions. Elle présente une singularité en $0$ mais comme c'est à l'intérieur du cylindre l'écoulement reste bien défini dans tout le domaine. 

Les composantes horizontales et verticales de la vitesse sont fournies par les parties réelles et imaginaires de la dérivée de $\vcomp(z)$:

\begin{equation}
  \begin{split}
    u + iv = U \left(1 - \frac{R^2}{z z}\right) &= U\left(1 - R^2\frac{ {z^\star}^2}{|z|^4}\right)
                               &= U (1 - R^2\frac{ (x^2 - y^2 - 2i xy )}{r^4} )
  \end{split}
\end{equation}

On trouve deux points de stagnation (vitesse nulle) sur l'axe $y=0$, pour $x\pm R$. 

\begin{figure}
\includegraphics{Python/03_zero_circulation.pdf}
\label{fig:zero_circulation}
\caption{Fonction de courant (traits pleins) et lignes de potentiel de vitesse (pontillés) pour un fluide parfait rencontrant un obstacle cylindrique, en l'absence de circulation autour de l'obstactle ($\vcomp(z) = U ( z + \frac{R^2}{z})$)}
\end{figure}

\todo[inline]{expliquer que c'est une transformation conforme. }



\section{Analyse de l'écoulement autour d'un cylindre avec circulation}

On peut étudier la position des points de stagnation dans ces conditions, comme illustré sur la figure~\ref{fig:cylindre_with_circulation}. Bien que l'analyse algébrique ne soit pas excessivement complexe, nous n'en fournisson ici que les conclusions. Pour $\Gamma > 0$, les deux points de stagnation, initialement situés de part et d'autre du cylindre lorsque $\Gamma = 0$, migrent progressivement vers le bas. Ils finissent par se rencontrer et se détachent du cylindre dès que $\Gamma$ dépasse $4\pi UR$.

\begin{figure}
  \begin{center}
    \includegraphics[width=0.95\textwidth]{Python/04_SF_with_circulation.pdf}
  \end{center}
  \caption{}\label{fig:cylindre_with_circulation}
\end{figure}


Bien qu'il y ait une circulation, l'écoulement, dans tout son domaine, demeure irrotationnel et incompressible. Il respecte ainsi toutes les conditions du théorème de Bernoulli généralisé. Nous savons donc qu'en tout point,

\begin{displaymath}
\gpot + \frac{p}{\rho} + \frac{u_h^2}{2} \quad \text{est constant,}
\end{displaymath}

où nous utilisons la notation $u_h$ pour désigner le vecteur $(u,v)$ à deux dimensions.

Supposons ici que $\gpot$ est constant, ainsi que $\rho$. Le carré de la vitesse $u_h^2$ nous est donné par $|\od{\vcomp}{z}|^2$, si bien que
\[
p - \rho \frac{1}{2}|\od{\vcomp}{z}|^2 = p_\infty - \rho \frac{U^2}{2}.
\]
Il est donc tout à fait possible de calculer la force de pression résultante (par unité de hauteur du cylindre) :
\[
\vec{F_r} = \oint p\, \vec{n} \, \dif s.
\]

\begin{redmargin}
La suite du calcul est assez technique. En particulier, elle mobilise des notions d'\cindex{analyse complexe} qui sont vues à la fin du cours LMAT1222. Nous esquissons seulement la démonstration, avec l'objectif de nous concentrer sur le résultat principal et qui nous intéresse ici: le théorème de Kutta-Zhukovsky. 
La méthode consiste à exprimer la force résultante d'abord comme un vecteur $(D,L)$, où $D$ représente la \cindex{traînée} (\emph{drag}) et $L$ la \cindex{portance} (\emph{lift}). On calcule le nombre complexe $D - iL$ en voyant que, lors de l'intégrale sur le contour :
\[
\dif D = -p \dif y \quad \text{et} \quad \dif L = p \dif x.
\]
On montre alors (référence à Kundu, p. 171) que
\[
D - iL = \oint p \, \dif z^\star = \frac{1}{2} \left(\od{\vpot}{z}\right)^2 \dif z.
\]
Cette première expression correspond au \cindex{théorème de Blasius}.

La dérivée $\left(\od{\vpot}{z}\right)$ est connue et vaut :

\begin{displaymath}
U - \frac{R^2}{z^2} + \frac{i\Gamma}{2\pi z}.
\end{displaymath}

L'intégrale fait intervenir cette dérivée au carré et se calcule par la méthode des résidus (référence à Arfken). Après calcul, on aboutit au \cindex{théorème de Kutta-Zhukovsky} :
\end{redmargin}

\begin{key}{Théorème de Kutta-Zhukovsky}
Pour un écoulement incompressible, irrotationnel (fluide parfait), autour d'un obstacle, avec circulation, les valeurs de traînée et de portance sont données par :
\[
D = 0, \quad L = \rho U \Gamma.
\]
\end{key}

On en déduit deux enseignements importants :
\begin{itemize}
    \item Dans cet écoulement parfait, il n'y a pas de traînée (pas de force horizontale). C'est le \cindex{paradoxe de d'Alembert} en mécanique des fluides.
    \item Sans circulation, il n'y a pas de portance : pas de force perpendiculaire à l'écoulement.
\end{itemize}

Ce théorème est très général et ne nécessite pas de supposer un obstacle cylindrique. Il est possible de calculer la portance pour des profils d'ailes obtenus par une technique connue sous le nom de \cindex{transformation conforme}. Ici, nous avons simplement résolu numériquement l'équation de Laplace pour obtenir un profil d'aile d'avion, mais le problème reste le même : sans circulation, il n'y a pas de portance.

\begin{figure}
    % --- Figure 1 : Source ---
    \begin{minipage}[t]{0.49\textwidth}
        \centering
        \includegraphics[width=\textwidth]{Python/02_ellipse_streamfunction.pdf}
        \caption{Résolution numérique de l'équation de Laplace autour d'un obstacle cylindrique, en supposant aucune circulation autour de l'obstacle. Observez les points de stagnation (intersection entre obstactle et obstactle) de part et d'autre de l'obstacle}
        \label{fig:ellipse}
    \end{minipage}%
    \hfill%
    %
    % --- Figure 2 : Circulation ---
    \begin{minipage}[t]{0.49\textwidth}
        \centering
        \includegraphics[width=\textwidth]{Python/05_airfoil_without_circulation.pdf}
        \caption{Même figure, mais pour un obstacle ayant la forme d'une aile d'avion. Ici encore, nous avons fait l'hypothèse de l'absence de circulation. La solution est irréaliste, car elle impose à l'écoulement de contourner la pointe de l'aile d'avion. L'hypothèse est que ce contournement sera, en pratique, rendu impossible par des effets de friction}
        \label{fig:airfoil_without_circulation}
    \end{minipage}
\end{figure}



Il faut donc créer une circulation. C'est l'intérêt de l'aspect anguleux à l'arrière de l'aile. On crée artificiellement un point de stagnation. La seule façon de satisfaire la condition du point de stagnation est de déplacer les points de stagnation, initialement situés de part et d'autre de l'aile, en créant une circulation. C'est cette circulation qui est nécessaire à la portance.

Mais comment induire cette circulation ? Dans l'\cindex{effet Magnus}, c'est le ballon qui cède du moment angulaire au fluide. Mais ici, il n'y a pas de force extérieure : l'aile ne tourne pas. La circulation est donc créée à l'interface aile-fluide par des frottements, mais doit être compensée (pour conserver le moment angulaire) par un contre-vortex qui se développe à l'arrière de l'aile.

\begin{figure}
    % --- Figure 1 : Source ---
  \scalebox{0.7}{\input{Pdftex/04_aile_avion.pdf_tex}}
   \caption{Schéma de l'écoulement attendu autour d'une aile d'avion. En bleu, nous avons représenté schématiquement la solution sans circulation. La circulation réelle sera plus proche de la solution noire, ou à la solution bleue s'ajoute une circulation (en rouge), telle que le point de stagnation arrière est déplacé sur la pointe de l'aile d'avion. Cette circulation nette explique la portance. }
   \label{fig:airfoil_schematic}
\end{figure}
